% Compile with pdflatex (run twice for the table of contents / \Cref).
\documentclass[11pt]{article}

% evan.sty by Evan Chen (Boost Software License 1.0)
% https://github.com/vEnhance/dotfiles/blob/main/texmf/tex/latex/evan/evan.sty
% Options: sexy = colored boxes for theorems + colored section heads,
%          hints, british, nothm, nofancy, ... (see the top of evan.sty)
\usepackage[sexy]{evan}

\title{Problem set 1 solution}
\author{Anakawee Chujinda}

% evan.sty hardcodes its author into the running header; reset it.
\lhead{Anakawee Chujinda}
\rhead{\nouppercase{\leftmark}}

\begin{document}
\maketitle
\tableofcontents
\newpage

\section{The real numbers}

\subsection{Ordered fields and completeness}

A quick reminder of notation: $\NN$, $\ZZ$, $\QQ$, $\RR$, $\CC$ are the usual
number systems, and $\eps > 0$ always denotes a real number.

\begin{definition}
  Let $S \subseteq \RR$ be nonempty. We say $b \in \RR$ is an
  \vocab{upper bound} of $S$ if $x \le b$ for every $x \in S$.
  The \vocab{supremum} $\sup S$ is the least upper bound of $S$,
  when one exists.
\end{definition}

\begin{example}
  [Rationals are not complete]
  The set $S = \{ q \in \QQ \mid q^2 < 2 \}$ is bounded above in $\QQ$
  but has no least upper bound in $\QQ$.
\end{example}

\begin{theorem}
  [Archimedean property]
  \label{thm:archimedean}
  If $x, y \in \RR$ and $x > 0$, then there exists $n \in \NN$ with $nx > y$.
\end{theorem}

\begin{proof}
  Suppose not, so $A = \{ nx \mid n \in \NN \}$ is bounded above by $y$.
  Let $\alpha = \sup A$. Since $x > 0$, the number $\alpha - x$ is not an
  upper bound, so $\alpha - x < mx$ for some $m \in \NN$. Then
  $\alpha < (m+1)x \in A$, contradicting $\alpha = \sup A$.
\end{proof}

\begin{corollary}
  $\QQ$ is dense in $\RR$: if $x < y$ then there is $q \in \QQ$ with
  $x < q < y$.
\end{corollary}

\begin{remark}
  \Cref{thm:archimedean} is exactly what fails in a non-Archimedean
  ordered field, e.g.\ the field of rational functions over $\RR$.
\end{remark}

\subsection{Sequences}

\begin{definition}
  A sequence $(a_n)$ \vocab{converges} to $L$ if for every $\eps > 0$
  there is $N$ such that $\abs{a_n - L} < \eps$ for all $n \ge N$.
\end{definition}

\begin{ques}
  Why can a sequence not converge to two different limits?
\end{ques}

\begin{exercise}
  Show that every convergent sequence in $\RR$ is bounded, and that the
  converse is false.
\end{exercise}

\begin{proposition}
  [Monotone convergence]
  Every bounded monotone sequence of reals converges.
\end{proposition}

\begin{proof}
  Say $(a_n)$ is increasing and bounded, and put $L = \sup_n a_n$.
  Given $\eps > 0$, some $a_N > L - \eps$; monotonicity gives
  \[ L - \eps < a_N \le a_n \le L \qquad \text{for all } n \ge N, \]
  hence $\abs{a_n - L} < \eps$.
\end{proof}

\section{Problem set}

\begin{problem}
  Let $f \colon \RR \to \RR$ be continuous with $f(x) = 0$ for all
  $x \in \QQ$. Prove $f \equiv 0$.
\end{problem}

\begin{soln}
  Fix $x \in \RR$ and take rationals $q_n \to x$. Continuity gives
  $f(x) = \lim_n f(q_n) = 0$.
\end{soln}

\end{document}
